% ── merged from midcurve_conclusion_ann.tex (Phase I summary) ──────────────────

%%%%%%%%%%%%%%%%%%%%%%%%%%%%%%%%%%%%%%%%%%%%%%%%%%%%%%%%%%%%%%%%%%%%%%%%%%%%%%%%%%
\begin{frame}[fragile]\frametitle{Phase I: Context}
	\begin{itemize}
	\item Various applications need lower dimensional representation of shapes.
	\item Midcurve is one-dimensional (1D) representation of a two-dimensional (2D) planar shape.
	\item Used in animation, shape matching, retrieval, finite element analysis, etc.
	\end{itemize}
\end{frame}

%%%%%%%%%%%%%%%%%%%%%%%%%%%%%%%%%%%%%%%%%%%%%%%%%%%%%%%%%%%%%%%%%%%%%%%%%%%%%%%%%%
\begin{frame}[fragile]\frametitle{Phase I: Approach}
	\begin{itemize}
	\item Prior approaches: Thinning, Medial Axis Transform (MAT), Chordal Axis Transform (CAT), Straight Skeletons --- all rule-based, not scalable.
	\item MidcurveNN: Encoder-Decoder neural network learns the mapping from 2D polygon image to 1D midcurve image in a supervised manner.
	\item Once trained, the network predicts midcurves for unseen polygon shapes.
	\end{itemize}
\end{frame}

%%%%%%%%%%%%%%%%%%%%%%%%%%%%%%%%%%%%%%%%%%%%%%%%%%%%%%%%%%%%%%%%%%%%%%%%%%%%%%%%%%
\begin{frame}[fragile]\frametitle{Phase I: Results}

\begin{center}
\small
\begin{tabular}{@{}lcccc@{}}
\toprule
\textbf{Variant} & \textbf{Input} & \textbf{Loss} & \textbf{IoU} & \textbf{Branches} \\
\midrule
Simple AE  & $100\times100$ & High   & Low  & \texttimes \\
Dense AE   & $100\times100$ & Medium & Med  & \texttimes \\
CNN AE     & $128\times128$ & Low    & High & \texttimes \\
\bottomrule
\end{tabular}
\end{center}

\begin{itemize}
\item CNN AE achieves $\approx10\times$ lower loss than Dense AE
\item All variants lose exact geometry (rasterisation trade-off)
\item None can represent branched midcurves (T, Plus) natively
\end{itemize}
\end{frame}

% ── merged from midcurve_conclusion_llm.tex (Phase II summary) ─────────────────

%%%%%%%%%%%%%%%%%%%%%%%%%%%%%%%%%%%%%%%%%%%%%%%%%%%%%%%%%%%%%%%%%%%%%%%%%%%%%%%%%%
\begin{frame}[fragile]\frametitle{Phase II: Motivation}
	\begin{itemize}
	\item Traditional methods are predominantly rule-based: not scalable to arbitrary shapes.
	\item A novel Large Language Model approach attempts to build a generic, data-driven model.
	\item GPT-4 (few-shot) succeeds on simple shapes; open-source models struggle without fine-tuning --- pointing to the value of domain-specific QLoRA fine-tuning.
	\end{itemize}
\end{frame}

%%%%%%%%%%%%%%%%%%%%%%%%%%%%%%%%%%%%%%%%%%%%%%%%%%%%%%%%%%%%%%%%%%%%%%%%%%%%%%%%%%
\begin{frame}[fragile]\frametitle{Phase II: Findings}
    \begin{itemize}
        \item BRep JSON representation solves the branching serialization problem for T- and Plus-shapes.
        \item QLoRA fine-tuning of Qwen2.5-7B achieves \textbf{PSR = 98\%, MAE = 0.78} on the 4-shape benchmark.
        \item Nemotron-Mini-4B (2-shot, no fine-tuning): PSR = 85.7\%, topology accuracy = 0.83, fits in $\sim$2\,GB VRAM.
        \item Key challenge: coordinates are text tokens --- digit arithmetic is harder for LLMs than neural regression.
    \end{itemize}
\end{frame}

%%%%%%%%%%%%%%%%%%%%%%%%%%%%%%%%%%%%%%%%%%%%%%%%%%%%%%%%%%%%%%%%%%%%%%%%%%%%%%%%%%
\begin{frame}[fragile]\frametitle{Phase II: MidcurveLLM Architecture}
    \begin{itemize}
        \item Utilizes an Encoder-Decoder architecture and B-rep structures.
        \item Showcases promise, despite discrepancies with ground truths.
    \end{itemize}
\end{frame}

%%%%%%%%%%%%%%%%%%%%%%%%%%%%%%%%%%%%%%%%%%%%%%%%%%%%%%%%%%%%%%%%%%%%%%%%%%%%%%%%%%
\begin{frame}[fragile]\frametitle{Phase II: Implications and Future Directions}
    \begin{itemize}
        \item Points to the need for more extensive datasets and refined training parameters.
        \item Serves as a crucial catalyst for advancing midcurve computation methodologies.
        \item Invites further scrutiny and advancements in the transformative intersection of geometry and advanced machine learning.
    \end{itemize}
\end{frame}

% ── overall conclusion (original midcurve_conclusion.tex) ──────────────────────

%%%%%%%%%%%%%%%%%%%%%%%%%%%%%%%%%%%%%%%%%%%%%%%%%%%%%%%%%%%%%%%%%%%%%%%%%%%%%%%%%%
\begin{frame}[fragile]\frametitle{Three Approaches: Summary}

\begin{center}
\small
\begin{tabular}{@{}p{1.8cm}p{2.2cm}p{2.0cm}p{2.2cm}@{}}
\toprule
 & \textbf{Phase I} & \textbf{Phase II} & \textbf{Phase III} \\
 & \textbf{Image} & \textbf{Text/LLM} & \textbf{Geometry} \\
\midrule
Input type      & Bitmap (100--256px) & BRep JSON text & Graph $(x,y)$ \\
Output type     & Bitmap             & BRep JSON text & Graph $(x,y)$ \\
Exact coords    & No                 & Yes            & Yes \\
Branches        & No                 & Yes (BRep)     & Yes (native) \\
Best model      & CNN encoder-decoder & Qwen2.5-7B QLoRA / Nemotron-Mini-4B & Graph Transformer \\
Sub-variants    & 7 variants         & Prompt, CodeT5, QLoRA, Nemotron3  & GT + Graphormer FT \\
Key limitation  & Rasterisation loss & Token arithmetic; repair needed & TopK fixes node count \\
Status          & Fully impl.        & Implemented    & Implemented \\
\bottomrule
\end{tabular}
\end{center}
\end{frame}

%%%%%%%%%%%%%%%%%%%%%%%%%%%%%%%%%%%%%%%%%%%%%%%%%%%%%%%%%%%%%%%%%%%%%%%%%%%%%%%%%%
\begin{frame}[fragile]\frametitle{Conclusions}
	\begin{itemize}
	\item Midcurve computation is a \textbf{dimension-reduction problem}: 2D $\rightarrow$ 1D, graph-to-graph
	\item \textbf{Phase I} (image): encoder-decoder with UNet achieves good results; rasterisation trade-off is the key limitation
	\item \textbf{Phase II} (LLM): BRep JSON representation works; GPT-4 solves simple shapes via prompting; Nemotron-Mini-4B achieves 85.7\% valid JSON + 0.83 topology accuracy base-model; QLoRA fine-tuning (Qwen, Gemma, Nemotron) shows promise but needs more data
	\item \textbf{Phase III} (geometry): Graph Transformer operates directly on polygon coordinates, handles branches natively, and outputs a clean polyline graph without post-processing
	\end{itemize}
\end{frame}

%%%%%%%%%%%%%%%%%%%%%%%%%%%%%%%%%%%%%%%%%%%%%%%%%%%%%%%%%%%%%%%%%%%%%%%%%%%%%%%%%%
\begin{frame}[fragile]\frametitle{Next Frontier: Graph-to-Graph Translation}

	\begin{itemize}
	\item The ultimate formulation is \textbf{graph-to-graph}: profile graph $\Rightarrow$ midcurve graph
	\item This is a \textbf{graph coarsening / dimension-reduction} problem, currently underexplored in geometric deep learning
	\item Promising directions:
		\begin{itemize}
		\item Non-auto-regressive Graph Transformers (Phase III, this work)
		\item Diffusion-based graph generation conditioned on input topology
		\item Geometric Graph Normalizing Flows
		\item Integration of LLM structural reasoning (Phase II) with geometric constraints (Phase III)
		\end{itemize}
	\item Goal: a model that generalises to \textbf{any 2D thin polygon} (not just I, L, T, Plus) and produces an exact, branch-aware midcurve without rasterisation or post-processing
	\end{itemize}

\begin{center}
\alert{\Large Graph-to-Graph is the next frontier for MidcurveNN.}
\end{center}
\end{frame}
